\begin{titlepage}
 
 
\setlength{\centeroffset}{-0.5\oddsidemargin}
\addtolength{\centeroffset}{0.5\evensidemargin}
\thispagestyle{empty}

\noindent\hspace*{\centeroffset}\begin{minipage}{\textwidth}

\centering


 \vspace{3.3cm}

%si el proyecto tiene logo poner aquí
\includegraphics{imagenes/logo.png} 
 \vspace{0.5cm}

% Title

{\Huge\bfseries Título del proyecto}
\noindent\rule[-1ex]{\textwidth}{3pt}\\[3.5ex]
{\large\bfseries Subtítulo del proyecto.\\[4cm]}
\end{minipage}

\vspace{2.5cm}
\noindent\hspace*{\centeroffset}\begin{minipage}{\textwidth}
\centering

\textbf{Autor}\\ {Ángel Gómez Martín}\\[2.5ex]
\textbf{Director}\\
{Juan Julián Merelo Guervos}\\[2cm]
%\includegraphics[width=0.15\textwidth]{imagenes/tstc.png}\\[0.1cm]
%\textsc{Departamento de Teoría de la Señal, Telemática y Comunicaciones}\\
%\textsc{---}\\
Granada, mes de 2020
\end{minipage}
%\addtolength{\textwidth}{\centeroffset}
\vspace{\stretch{2}}

 
\end{titlepage}




\cleardoublepage

\chapter*{}
\thispagestyle{empty}

\begin{center}
{\large\bfseries IPManager: Infraestructura bajo demanda}\\
\end{center}
\begin{center}
Ángel Gómez Martín
\end{center}

\noindent{\textbf{Palabras clave}: servicios, infraestructura, aprovisionamiento, backend, frontend, \textit{API REST}}\\

\vspace{0.7cm}
\noindent{\textbf{Resumen}}\\

El aprovisionamiento, despliegue de servicios y almacenamiento de configuraciones de máquinas a gran escala son tareas que tradicionalmente se han desarorollado de forma independiente. Esto supone a corto y largo plazo una gran inversión de tiempo para los administradores de sistemas.

\bigskip
En este proyecto se ha desarrollado una solución compuesta por un backend y un frontend. Esta busca unificar todos estos procesos en un puesto centralizado que permite agilizar el desarrollo de estas tareas. También ofrece modularidad en dos grandes aspectos: permitiendo que se puedan agregar nuevas funcionalidades de forma sencilla en el backend y proporcionando una \textit{API REST} que permite la comunicación con otro tipo de aplicaciones.


\cleardoublepage
\thispagestyle{empty}

\begin{center}
{\large\bfseries Project Title: Project Subtitle}\\
\end{center}
\begin{center}
Ángel Gómez Martín
\end{center}

\noindent{\textbf{Keywords}: services, infrastructure, provisioning, backend, frontend, \textit{REST API}}\\

\vspace{0.7cm}
\noindent{\textbf{Abstract}}\\

Provisioning, deployment of services and storage of large-scale machine configurations are tasks that have traditionally been carried out independently. In the short and long term, this is a major time investment for system administrators.

\bigskip
In this project, a solution composed of a backend and a frontend has been developed. It seeks to unify all these processes in a centralized position that allows to speed up the course of these tasks. It also offers modularity in two major aspects: allowing new functionalities to be added easily in the backend and providing a \textit{REST API} that allows communication with other type of applications.











\chapter*{}
\thispagestyle{empty}

\noindent\rule[-1ex]{\textwidth}{2pt}\\[4.5ex]

Yo, \textbf{Ángel Gómez Martín}, alumno de la titulación Grado en Ingeniería Informática de la \textbf{Escuela Técnica Superior de Ingenierías Informática y de Telecomunicación de la Universidad de Granada}, con DNI XXXXXXXXX, autorizo la ubicación de la siguiente copia de mi Trabajo Fin de Grado en la biblioteca del centro para que pueda ser consultada por las personas que lo deseen.

\vspace{6cm}

\noindent {Fdo: Ángel Gómez Martín}

\vspace{2cm}












\chapter*{}
\thispagestyle{empty}

\noindent\rule[-1ex]{\textwidth}{2pt}\\[4.5ex]

D. \textbf{Juan Julián Merelo Guervós}, Profesor del Área de XXXX del Departamento de Arquitectura y Tecnología de Computadores de la Universidad de Granada.

\vspace{0.5cm}

\textbf{Informa:}

\vspace{0.5cm}

Que el presente trabajo, titulado \textit{\textbf{IPManager, Infraestructura bajo demanda}},
ha sido realizado bajo su supervisión por \textbf{Ángel Gómez Martín}, y autoriza la defensa de dicho trabajo ante el tribunal que corresponda.

\vspace{0.5cm}

Y para que conste, expide y firma el presente informe en Granada a X de Julio de 2020.

\vspace{1cm}

\textbf{El director:}

\vspace{5cm}

\noindent {Fdo: Juan Julián Merelo Guervós}








\chapter*{Agradecimientos}
\thispagestyle{empty}
\vspace{1cm}

A mi tutor, por la ayuda, conocimientos e ideas prestadas para la realización de este proyecto.

\bigskip
A mis padres y mi hermana, por todo el cariño y apoyo que me han dado durante toda esta etapa.

